\section{Integer signatures and difference counts}
\label{sec:signatures}

The geometric comparison will require log-concavity of lattice-pair
counts along the directions $e_j-e_k$.  The eventual counting algorithm
will require two quadratic inequalities for positive weights on integer
boxes.  We develop a common argument for both purposes: rank penalties
satisfy the inequalities, and summing paired coordinates preserves them.
The quadratic viewpoint comes from the rank-deficiency construction
in~\cite[Section~3, ``Quadratic signature inequalities'']{Parent}.
That paper derives its forms from the Lorentzian property of homogenized
matroid Tutte polynomials and closure under nonnegative linear
substitutions~\cite[Theorems~4.10 and~2.10]{BH2020}.
Here we prove the integer statements directly.

\subsection{The two tests and balanced summation}

For a finite index set $I$, capacities $m\in\Z_{\ge0}^I$, and an
integer $r$, write
\[
 \mathcal B(m,r)=\{z\in\Z^I:0\le z_i\le m_i,\ \sum_i z_i=r\}.
\]
A \emph{box-valid display} is any integer vector in
$\prod_i[0,m_i]$, with its total specified separately.
Weights on $\mathcal B(m,r)$ are extended by zero outside this slice.

\begin{definition}[Integer quadratic signatures]
\label{def:integer-signatures}
A strictly positive weight $F$ on all of $\mathcal B(m,r)$ has the
\emph{removal signature} if, for every box-valid display $z$ of total
$r+2$, the matrix
\begin{equation}
 \bigl(F(z-e_i-e_j)\bigr)_{i,j\in I}
 \quad\text{has at most one positive eigenvalue}.
 \label{eq:integer-removal-test}
\end{equation}
It has the \emph{addition signature} if, for every box-valid display
$z$ of total $r-2$, the same conclusion holds for
\begin{equation}
 \bigl(F(z+e_i+e_j)\bigr)_{i,j\in I}.
 \label{eq:integer-addition-test}
\end{equation}
Diagonal entries represent two changes in the same coordinate.
All invalid changes have weight zero.  A test with no box-valid
display is vacuous.
\end{definition}

If $Q$ is symmetric with at most one positive eigenvalue and
$u^{\mathsf T}Qu>0$, then $Q$ is nonpositive on
$\{v:v^{\mathsf T}Qu=0\}$.  Otherwise this subspace and $u$ would
contain a two-dimensional positive-definite subspace.  Consequently
\begin{equation}
 v^{\mathsf T}Qv\le
 \frac{(v^{\mathsf T}Qu)^2}{u^{\mathsf T}Qu}.
 \label{eq:signature-projection}
\end{equation}
We will also use its null-axis version: if
$e^{\mathsf T}Qe=0$ and $Qe\ne0$, then $Q$ is nonpositive on
$\{v:v^{\mathsf T}Qe=0\}$.  Indeed, take $w=Qe$ and apply
\eqref{eq:signature-projection} with $u=e+tw$ for sufficiently
small $t>0$.  For $v^{\mathsf T}Qe=0$, its right-hand side is
$t^2(v^{\mathsf T}Qw)^2/(2t\|Qe\|^2+t^2w^{\mathsf T}Qw)$,
which tends to zero.

For two coordinates $a,b$ of the same capacity $M$, summing
\emph{at balance} means summing their values $(k,M-k)$,
$0\le k\le M$.  The next lemma separates this operation from the
particular factors used to construct the weight.
Its proof is the balanced-summation induction
of~\cite[Section~3, Lemma~3.1, ``Integer signature under balanced
summation'']{CB}, applied to a general positive weight.
Only positivity and the initial quadratic signature enter the induction.

\begin{lemma}[Balanced summation]
\label{lem:balanced-summation}
Let $F>0$ on the entire slice $\mathcal B(m,r)$ and suppose it has
the removal signature.  Choose disjoint coordinate pairs
$\{a_t,b_t\}$ with $m_{a_t}=m_{b_t}=M_t\ge1$.
Their balanced sum
\[
 \Phi(z)=\sum_{0\le k_t\le M_t}
 F\bigl(z,(k_t,M_t-k_t)_t\bigr)
\]
is strictly positive on the entire remaining box slice of total
$r-\sum_tM_t$ and has the removal signature.
The same statement holds with ``addition'' in place of ``removal''.
Fixing coordinates at prescribed box-valid values preserves either
signature and positivity on the resulting valid slice.
\end{lemma}

\begin{proof}
Fixing coordinates restricts each test to a principal submatrix of
the original test.  For summation it suffices to remove one pair at
a time.  Suppose earlier pairs have been summed, giving a positive
weight $\Psi$, and the next pair has capacity $M$.
Fix a box-valid display $z$ on the other coordinates, two above
their required total after this pair is summed.
For $0\le k\le M$, the removal matrix at $(z,k,M-k)$ is
\begin{equation}
 N_k=
 \begin{pmatrix}
  V_k&S(k-1)&S(k)\\
  S(k-1)^{\mathsf T}&u_{k-1}&u_k\\
  S(k)^{\mathsf T}&u_k&u_{k+1}
 \end{pmatrix},
 \label{eq:balanced-matrices}
\end{equation}
where
\begin{align*}
 (V_k)_{ij}&=\Psi(z-e_i-e_j,k,M-k),\\
 S(l)_i&=\Psi(z-e_i,l,M-1-l),\\
 u_k&=\Psi(z,k-1,M-1-k).
\end{align*}
Each $N_k$ has at most one positive eigenvalue by induction.
Out-of-box terms are zero, so $S(-1)=S(M)=0$ and $u_k=0$
unless $1\le k\le M-1$.
The matrix we must control is $\sum_{k=0}^M V_k$.
We will find one common subspace of codimension at most one on
which its quadratic form is nonpositive.

When $M=1$, put $S=S(0)$.  After discarding a zero axis, the two
matrices in~\eqref{eq:balanced-matrices} are
\[
 \begin{pmatrix}V_0&S\\S^{\mathsf T}&0\end{pmatrix},
 \qquad
 \begin{pmatrix}V_1&S\\S^{\mathsf T}&0\end{pmatrix}.
\]
If $S\ne0$, the null-axis form of
\eqref{eq:signature-projection} makes both $V_0$ and $V_1$
nonpositive on $\{a:a^{\mathsf T}S=0\}$.
If some $z_i>0$, positivity gives
$S_i=\Psi(z-e_i,0,0)>0$.  Thus $S=0$ implies $z=0$, and then
both $V_k$ are zero because all removals are invalid.
This proves the assertion for $M=1$.

Let $M\ge2$.  Positivity on the full slice gives $u_k>0$ for
$1\le k\le M-1$.  Take a vector $a$ on the old coordinates with
\[
 a^{\mathsf T}\sum_{l=0}^{M-1}S(l)=0,
\]
and define
\[
 g_k=\frac{a^{\mathsf T}\sum_{l=0}^{k-1}S(l)}{u_k}
 \quad(1\le k\le M-1),
 \qquad g_k=0\quad\text{otherwise}.
\]
The definition and the condition on $a$ give, including the endpoints,
\begin{equation}
 a^{\mathsf T}S(l)=u_{l+1}g_{l+1}-u_lg_l
 \qquad(0\le l\le M-1).
 \label{eq:balanced-cumulative}
\end{equation}
For $v_k=(a,g_k,-g_k)$, its products under the form $N_k$ with
the two new coordinate axes are
\[
 u_{k-1}(g_k-g_{k-1}),\qquad
 u_{k+1}(g_{k+1}-g_k),
\]
respectively.  If at least one axis has positive square,
\eqref{eq:signature-projection} applied to that axis bounds
$v_k^{\mathsf T}N_kv_k$ by its corresponding term in
\begin{equation}
 v_k^{\mathsf T}N_kv_k
 \le u_{k-1}(g_k-g_{k-1})^2+
      u_{k+1}(g_k-g_{k+1})^2.
 \label{eq:balanced-projection}
\end{equation}
Adding the other, nonnegative term preserves the bound.
If both axis squares vanish, $u_{k-1}=u_{k+1}=0$.
At $k=0$ or $M$ the other neighboring $u$ is positive, so this
case is interior and $u_k>0$.  The sum of the two axes then has
square $2u_k>0$ and product zero with $v_k$.
Applying~\eqref{eq:signature-projection} to that sum proves
\eqref{eq:balanced-projection}, with right-hand side zero.

Expanding~\eqref{eq:balanced-projection} and substituting
\eqref{eq:balanced-cumulative} cancels the mixed terms and gives
\[
 a^{\mathsf T}V_ka+2u_kg_k^2
   -u_{k-1}g_{k-1}^2-u_{k+1}g_{k+1}^2\le0.
\]
Summing over $0\le k\le M$ cancels every $u_kg_k^2$ term.
Hence $\sum_kV_k$ is nonpositive on the kernel of the single
linear functional $a\mapsto a^{\mathsf T}\sum_lS(l)$, as required.
Every remaining box-valid profile has positive weight because any
balanced choices of the summed pairs supply a positive summand.
This completes the removal induction.

For addition, reflect all coordinates in their boxes,
$z_i\mapsto m_i-z_i$.  This interchanges the two signatures.
It preserves balance and commutes with the balanced sum, so the
removal result for the reflected weight gives the addition result.
\end{proof}

The one-positive-eigenvalue property is preserved by matrix limits.
All sums above are finite, so we may also pass to limits of the
positive weights after summation.  Positivity is needed \emph{before}
this passage: the proof divides by the interior quantities $u_k$.
We will use global interpolation for positive approximations before
taking limits, since adjacent log-concavity inequalities alone would
not rule out gaps in the limiting support.

\subsection{Rank penalties, grouped coordinates, and reflection}

We next construct positive weights to which the summation lemma
applies.  Let $r$ be an integral normalized nondecreasing submodular
function on a finite ground set $E$.  For $v\in\Z_{\ge0}^E$, put
\[
 m_r(v)=\min_{A\subseteq E}\{r(A)+v(E\setminus A)\}.
\]
By Lemma~\ref{lem:expansion-matroid}, this is the rank of a set
with occupancy vector $v$ in a sufficiently large copy matroid.
For $H\ge0$, define
\begin{equation}
 \Phi_{r,H}(v)=
 \exp\!\left\{-H\left[2\bigl(v(E)-m_r(v)\bigr)
                         +r(E)-v(E)\right]\right\}.
 \label{eq:rank-soft-weight}
\end{equation}
The bracket equals $(v(E)-m_r(v))+(r(E)-m_r(v))$ and is
nonnegative.  It vanishes exactly on $B(r)\cap\Z^E$.
Thus $\Phi_{r,H}$ is positive for finite $H$ and converges to the
indicator of integer bases as $H\to\infty$.

For a positive factor $G$ defined on a whole box, the normalized
matrix for steps of sign $\sigma\in\{1,-1\}$ is
\begin{equation}
 Q_{ij}=\frac{G(z)G(z+\sigma e_i+\sigma e_j)}
                   {G(z+\sigma e_i)G(z+\sigma e_j)}.
 \label{eq:signature-normalization}
\end{equation}
Only coordinates permitting one step are retained; an invalid
repeated step is assigned $Q_{ii}=0$.
The unnormalized test is obtained from $Q$ by a positive scalar
and positive diagonal congruence.  In particular,
$\one\one^{\mathsf T}-Q\succeq0$ implies the required signature.

\begin{lemma}[Rank-factor signatures]
\label{lem:rank-signatures}
Partition a set of count coordinates $I$ into groups indexed by $E$,
allowing empty groups, and let $b\in\Z_{\ge0}^E$.
For either of the two consistent orientations
\[
 v_e=b_e+\sum_{i\text{ in group }e}z_i,
 \qquad\text{or}\qquad
 v_e=b_e+\sum_{i\text{ in group }e}(m_i-z_i),
\]
the factor $G(z)=\Phi_{r,H}(v)$ satisfies
$\one\one^{\mathsf T}-Q\succeq0$ for both signs in
\eqref{eq:signature-normalization}.
The product of such factors on disjoint sets of count coordinates
has both integer signatures on every total slice.  The assertion
holds for every nonnegative power of the product, and for any
coordinates fixed at box-valid values.

More generally, a base constraint for an arbitrary integral
normalized submodular rank, with a fixed integer translation and
either consistent orientation of its grouped arguments on a finite
box, can be imposed as a limit of factors of this kind.
\end{lemma}

\begin{proof}
First take the increasing orientation and an addition test.
Realize the displayed $v$ by a set $S$ in the expansion matroid,
and assign two distinct spare copies to each retained count coordinate,
with these assignments disjoint even when several count coordinates
feed the same physical coordinate.
Contract $S$; equivalently, contract a basis of $S$ and delete its
remaining loops.  A new copy increases rank exactly when it is a
nonloop in this contraction.  Two individually rank-increasing
copies increase rank by only one precisely when they are parallel,
meaning that their two-element set has rank one.
In every other case their rank increments add.

The affine terms in the exponent of~\eqref{eq:rank-soft-weight}
cancel in~\eqref{eq:signature-normalization}.  Thus, when both
steps are valid, $Q_{ij}$ is $\exp(-2H)$ if their two nonloop
copies are parallel, and is $1$ otherwise.  For a repeated step
we use two distinct copies from that count coordinate.
Copies from one physical coordinate are interchangeable by matroid
automorphisms fixing $S$.  If the representative of one count coordinate
is parallel to that of another, exchange the first representative
with its second copy while fixing the other representative.
This shows that the second copy is parallel to the other representative
and hence to the first.  It follows that each nontrivial parallel class of retained
count indices has normalized entry $\exp(-2H)$ throughout its
block, including every valid repeated-step diagonal.
A single index whose two copies are parallel contributes a
singleton block.  All remaining valid diagonal entries are $1$.
Consequently
\begin{equation}
 \one\one^{\mathsf T}-Q
 =(1-e^{-2H})\sum_{C}\one_C\one_C^{\mathsf T}
       +\diag(\eta_i),\qquad \eta_i\ge0.
 \label{eq:rank-parallel-psd}
\end{equation}
Here $C$ ranges over these disjoint parallel classes.
The diagonal correction replaces an invalid repeated-step entry
by zero; this only increases the diagonal of the left-hand side.
This proves the addition assertion, including box boundaries.

For removals, use a finite expansion on a ground set $\Omega$
with enough copies to realize every display and with rank $r(E)$.
The complement $\Omega\setminus S$ in its dual matroid has rank
\[
 m^*(\Omega\setminus S)=|\Omega\setminus S|-r(E)+m_r(v).
\]
Writing $r^*=|\Omega|-r(E)$, we obtain the exact identity
\[
 |S|+r(E)-2m_r(v)
 =|\Omega\setminus S|+r^*-2m^*(\Omega\setminus S).
\]
A removal from $S$ is an addition to its dual complement.
Realize $S$ using $b_e$ fixed copies and disjoint sets of $z_i$
copies for the count coordinates in physical group $e$.
Choose one representative from every nonempty such set, and a second
copy when $z_i\ge2$.  A transposition of those two copies fixes
every other representative, preserves $S$, and preserves the dual
contraction.  The preceding parallel-class argument therefore gives
the same block description~\eqref{eq:rank-parallel-psd}.
When only one removal is allowed in a count coordinate, its diagonal
is zero and the nonnegative correction in that formula handles it.
No nonexistent second copy is required.

Reflection $z_i\mapsto m_i-z_i$ on the entire coordinate set
of one factor exchanges addition and removal, proving both tests
for the decreasing orientation.  Fixed coordinates merely become
part of $b$.  For factors on disjoint groups, normalized entries
between different groups equal one; hence the matrix
$\one\one^{\mathsf T}-Q$ is the direct sum of their nonnegative
definite matrices.  Unused coordinates contribute zero blocks,
apart from the permitted diagonal corrections.
This proves the product assertion.  Raising to a nonnegative
power replaces each $H$ by that power times $H$.

Finally let $f$ be an integral normalized submodular function that
need not be nondecreasing.  Choose an integer vector $a$ large
enough that $f'(A)=f(A)+a(A)$ is nondecreasing and that every
translated lower endpoint of the grouped arguments plus $a$ is
nonnegative.  This is possible: there are finitely many marginal
differences $f(A\cup\{e\})-f(A)$ and finitely many box endpoints.
The equation $B(f')=B(f)+a$ follows from
Lemma~\ref{lem:base-calculus}.  Apply the proved construction to
the shifted argument and $f'$, then let $H\to\infty$.
This gives exactly the desired base indicator on the finite box.
The copy matroids and these possibly large boxes are used only
to prove the inequalities; the algorithm does not construct them.
\end{proof}

\subsection{Log-concavity of difference counts}

We now combine the preceding lemmas.  The result concerns ordinary
counts of integer vectors, with no copy multiplicities.

\begin{proposition}[Root-direction log-concavity]
\label{prop:shift-logconcavity}
Let $f_1,f_2$ be integral normalized submodular functions on $E$,
not necessarily nondecreasing, and define
\[
 Y(d)=\#\{(x_1,x_2)\in(B(f_1)\cap\Z^E)\times
                         (B(f_2)\cap\Z^E):x_1-x_2=d\}.
\]
For $d_0\in\Z^E$ and distinct $j,k\in E$, put
$y(t)=Y(d_0+t(e_j-e_k))$, $t\in\Z$.
The positive support of $y$ is an integer interval, possibly
empty or a singleton.  For integers $a<t<b$ with
$y(a),y(b)>0$,
\begin{equation}
 y(t)\ge y(a)^{(b-t)/(b-a)}y(b)^{(t-a)/(b-a)}.
 \label{eq:shift-interpolation}
\end{equation}
\end{proposition}

\begin{proof}
Every base polytope here is bounded: its $e$th coordinate lies
between $f_i(E)-f_i(E\setminus\{e\})$ and $f_i(\{e\})$.
Fix any finite integer interval $[A,B]$ of parameter values,
with $B-A=M\ge2$.  For each $e\in E$, choose an integer interval
$[L_e,U_e]$ containing all $e$th coordinates of $B(f_2)$ and
having width $D_e=U_e-L_e\ge2$.
Introduce two count coordinates $a_e,b_e$, each of capacity $D_e$,
and two further coordinates $p,q$, each of capacity $M$.
Use the full box slice of total $\sum_eD_e+M$.
Define physical arguments by
\begin{align}
 w_{2,e}&=U_e-b_e,\notag\\
 w_{1,e}&=L_e+a_e+(d_0)_e
             +A\one_{\{e=j\}}-B\one_{\{e=k\}}
             +p\one_{\{e=j\}}+q\one_{\{e=k\}}.
 \label{eq:difference-pair-encoding}
\end{align}
The first rank factor uses the increasing counts $(a_e)_e,p,q$;
the second uses the decreasing counts $(b_e)_e$.
These sets of count coordinates are disjoint.
In the first factor, the coordinates $a_j,p$ feed the same physical
coordinate, as do $a_k,q$, exactly as allowed in
Lemma~\ref{lem:rank-signatures}.

On this finite box, choose modular translations $c_1,c_2$ so that
$f_i+c_i$ is nondecreasing and every lower endpoint of $w_i+c_i$
is nonnegative.  For finite $H$, the product
\[
 F_H=\Phi_{f_1+c_1,H}(w_1+c_1)
             \Phi_{f_2+c_2,H}(w_2+c_2)
\]
is strictly positive on the full slice and has both signatures.
Sum every ordinary pair $(a_e,b_e)$ at balance, leaving the pair
$(p,q)$ of total $M$.  By Lemma~\ref{lem:balanced-summation},
the resulting weight has the removal signature.
Let $y_H(t)$ denote its value at $(p,q)=(t-A,B-t)$.
It is positive for every integer $A\le t\le B$.

For $A<t<B$, add one to each coordinate of this retained pair.
This is a box-valid display of total $M+2$, whose removal matrix is
\[
 \begin{pmatrix}
  y_H(t-1)&y_H(t)\\y_H(t)&y_H(t+1)
 \end{pmatrix}.
\]
It has positive trace and at most one positive eigenvalue, so
$y_H(t)^2\ge y_H(t-1)y_H(t+1)$.
Because all entries are positive, the successive differences of
$\log y_H$ are nonincreasing.  It follows that, for all
$A\le a<t<b\le B$,
\[
 y_H(t)\ge y_H(a)^{(b-t)/(b-a)}
                 y_H(b)^{(t-a)/(b-a)}.
\]

At balance $a_e+b_e=D_e$, the vector $w_2$ equals $L+a$.
At $(p,q)=(t-A,B-t)$, formula~\eqref{eq:difference-pair-encoding}
therefore gives $w_1-w_2=d_0+t(e_j-e_k)$.
Each possible integer $x_2$ in the chosen intervals has exactly
one balanced encoding.  Letting $H\to\infty$ thus gives
$y_H(t)\to y(t)$ term by term, since the finite product of rank
penalties converges to the two base indicators.
Pass the global interpolation inequality to this limit.
Positive endpoint counts give a positive lower bound at every
intermediate integer.  The interval $[A,B]$ was arbitrary, proving
both~\eqref{eq:shift-interpolation} and the assertion about support.
\end{proof}

Proposition~\ref{prop:shift-logconcavity} supplies both the interval
support and the secant inequalities used in Section~\ref{sec:geometry}
to compare nearby shift counts. The positive approximations are what
make both conclusions survive the limit. The integrated weights of
Section~\ref{sec:ports} require a separate quantitative argument for
their signatures.
